% ══════════════════════════════════════════════════════════════════════════════
\chapter{System Design}
\label{ch:system-design}
% ══════════════════════════════════════════════════════════════════════════════

\section{System Architecture}

The station is built around a Raspberry Pi 3A+ that bridges three independent subsystems: a 2-axis AZ/EL rotator that aims a directional antenna at the satellite, an RF front-end that demodulates the signal received through that antenna, and a Python software stack that decides where to point and what to listen for. Figure~\ref{fig:system-arch} shows the data and power flow across the whole station; the rest of this section walks through it block by block.

\begin{figure}[H]
    \centering
    \includegraphics[width=\textwidth]{figures/system_architecture.png}
    \caption{Station architecture: Pi 3A+ host, MotorPCB, dual-CC1200 RF HAT, and power tree.}
    \label{fig:system-arch}
\end{figure}

\paragraph{Antenna chain.}
The AetherSpace CubeSat radiates a 433~MHz CC1200 downlink which, after $\sim$1000~km of free-space path, arrives at a directional UHF Yagi (Siretta Oscar 44, 9~dBi) bolted to the elevation cradle of the rotator. The Yagi feeds the RF HAT through a 3~m RG58 coax run terminated at an SMA connector on the HAT. The dashed link from the satellite to the antenna in Figure~\ref{fig:system-arch} marks this path as the link budget of \S\ref{sec:link-budget}.

\paragraph{RF HAT.}
A Raspberry Pi HAT carries the receive electronics: a CC1200 sub-GHz transceiver on a controlled-impedance RF section (with a second CC1200 footprint routed for VHF but not populated on the v1 board), a 5~V LMR51420 buck converter, a Raspberry Pi Pico (RP2040) acting as the radio's local controller, and two operator-feedback peripherals on the RP2040 GPIO (a piezo buzzer for AOS/LOS/error patterns and four WS2812B NeoPixels for status colour). The CC1200 talks to the RP2040 over a 5~MHz SPI1 bus; the RP2040 in turn frames a binary COBS protocol with CRC-16/CCITT-FALSE and pushes it to the Pi through the dedicated UART on \file{/dev/serial0}. Chapter~\ref{ch:rf-frontend} covers the RF schematic, layout, and matching network in detail.

\paragraph{Motor controller PCB.}
A second RP2040 (also on a Pi Pico module) runs the closed-loop tracking firmware: it reads two FAPG36-555-EN brushed DC gearmotors over their built-in quadrature encoders (64~CPR after 4$\times$ decode), drives them through two TB6642FG H-bridges at 20~kHz PWM, and exposes the standard EasyComm~III serial protocol to the Pi through USB CDC. An external ADXL345 IMU breakout provides the gravity vector for elevation homing at boot (SPI0). The 24~V supply that feeds the H-bridges is also stepped down on-board by a TPS560430 5~V/2~A buck for the Pico's logic rail. The two-axis output drives the AZ ($\pm$360\degree, 15:40 external reduction, 244.6 ticks/\degree) and EL (0--180\degree, 40:55 external reduction, 126.1 ticks/\degree) gear trains of the 3D-printed rotator described in \S\ref{sec:motors} and \S\ref{sec:gearing}.

\paragraph{Station computer.}
The Raspberry Pi 3A+ runs Debian Trixie (armhf, 512~MB) and hosts the application stack. \code{rotctld} from the hamlib package translates the rotator-side TCP control protocol on port~4533 into EasyComm~III bytes on the USB CDC link to the motor PCB. \code{station.py}, the tracking daemon, owns the orbital state machine, runs PyEphem at 10~Hz to compute commanded AZ/EL and Doppler-corrected frequency, drives \code{rotctld} and the RF HAT in lockstep, and emits a 1~Hz status snapshot at \file{\textasciitilde/.station\_status.json}. \code{dashboard.py} reads that snapshot and serves a self-signed HTTPS dashboard on port~5000 to a phone or laptop browser. \code{satnogs.py} pushes CRC-valid frames to the SatNOGS DB (station ID 4712) over the SiDS HTTP API, with an offline queue for sessions without internet. PyEphem is fed by TLE data periodically refreshed from AMSAT and CelesTrak, and per-satellite radio profiles come from a hand-curated \file{sat\_profiles.json} with a SatNOGS DB transmitter-API fallback. Chapter~\ref{ch:firmware-software} expands the Pi stack.

\paragraph{External services and operator.}
Three external hosts close the loop. AMSAT and CelesTrak supply orbital elements (TLEs) over plain HTTP; SatNOGS DB receives decoded telemetry over HTTPS. The operator's phone or laptop browser connects to the Pi's HTTPS dashboard for setup, GPS entry, and live monitoring; SSH is not used in the deployment workflow.

\paragraph{Power tree.}
A single 24~V bench supply (or 6S LiPo for portable operation) feeds the motor PCB's 24~V H-bridge rail directly and is dropped to 5~V on each board independently: an LMR51420 on the RF HAT and a TPS560430 on the motor PCB. The 5~V rail of the RF HAT also supplies the Pi through the GPIO header. Splitting the buck converters this way prevents the high-current H-bridge switching transients on the motor PCB from propagating into the RF HAT's 5~V rail; each board's switching regulator is local to its own load.

\paragraph{Boot order.}
Five systemd units start the station autonomously after the Pi finishes booting: \code{rfkill-unblock} clears the default-blocked WiFi on Trixie, \code{wifi-provision} hosts a captive-portal SSID/password entry if no WiFi is configured, \code{rotctld} starts the hamlib daemon, \code{dashboard} starts the HTTPS server, and \code{station} starts the tracking daemon once \file{station.conf} (written by the dashboard's setup wizard from the operator's phone GPS) is available.

\paragraph{Subsystem independence.}
The architecture deliberately keeps rotator and RF independent: each subsystem has its own RP2040 and its own protocol to the Pi (USB CDC EasyComm vs.\ UART COBS), so either side can be tested, replaced, or upgraded in isolation. The v2 MOTOR\_RF\_HAT collapses both PCBs onto a single 6-layer Pi HAT (\S\ref{sec:cost}), but the firmware boundary stays the same: the merged board still presents two logically distinct devices to the Pi.

\paragraph{Design constraints.}
Four constraints shape the architecture. Portability is enforced by a 1/4"-20 UNC tripod thread on the base, with the full setup measured at approximately 4 minutes from unpacking to first tracked pass across three field sessions (\S\ref{sec:field-tests}). The zero-SSH workflow routes every field operation, GPS entry, north alignment, pass selection, park, and shutdown, through the HTTPS dashboard, which targets a laptop browser but also works on a phone for quick tasks such as supplying a GPS fix. The cost target is the v2 MOTOR\_RF\_HAT BOM derived in \S\ref{sec:cost} and listed in Appendix~\ref{app:bom}: $\sim$€40 for the combined controller/RF PCB, $\sim$€75 for the rotator hardware, $\sim$€100 for the station core (excluding antenna, coax, power supply, tripod). The intended licensing is CERN-OHL-S v2 for hardware and MIT for firmware and Python stack, with the formal \file{LICENSE} files added on the open-source release.

\section{Mechanical Design}

The rotator is designed in Onshape and printable on consumer FDM machines with a 220$\times$220~mm bed (the largest single printed part is 205$\times$160$\times$120~mm). All structural prints are in ASA (Acrylonitrile Styrene Acrylate), an outdoor-grade engineering filament for desktop FDM. ASA is chosen over PLA for two properties: (i) UV stability from the acrylate rubber phase (a polymer-science property of ASA copolymers; ASA substitutes the acrylate rubber for the butadiene phase of ABS, eliminating the UV-degradation vector), and (ii) a glass transition $T_g \approx 100$\degC / Vicat softening $\approx$~105\degC, compared to $T_g \approx 60$\degC\ for PLA; this margin avoids thermal deformation in direct sunlight. PETG was considered and rejected because of lower interlayer adhesion and creep under sustained load compared with ASA at equivalent print settings.

\begin{figure}[H]
    \centering
    \includegraphics[width=0.7\textwidth]{photos/rotator_hero_2.jpg}
    \caption{Assembled 3D-printed rotator, front view.}
\end{figure}

\subsection{Structure}

The chassis consists of a base plate with an integrated azimuth bearing race, two side plates forming the elevation frame, and a removable hood covering the electronics compartment. The hood attaches via a 3D-printed hinge, giving tool-free access to cable connections during bring-up.

The station mounts to any standard camera tripod through a 1/4"-20 UNC brass heat-set insert in the base plate. Using the tripod standard rather than a custom mount lets the station share hardware with existing camera tripods in the user's toolkit.

The elevation axis is a PVC pipe passing through both side plates. The elevation gear is clamped directly onto this pipe by integrated 3D-printed screw clamps. The clamp geometry lets the EL gear be slid along the pipe during assembly to set mesh with the motor pinion, without set screws or adhesives.

All fasteners are M4$\times$10~mm countersunk screws for structural joints and M3$\times$8~mm for motor mounts. The assembly uses no adhesives (every joint is either a press-fit, a screw fastener, or a clamp), so the station disassembles for transport or repair without destructive separation.

\subsection{Bearings}

Both axes use 3D-printed deep-groove ball bearings with COTS 3~mm steel balls. The elevation outer race is printed in place into the side plate. The azimuth outer race is a separate screw-in part: the raceway has to be printed flat-down on the build plate to avoid layer stairstepping on the running surface, which prevents printing it as a feature of the base plate. Inner races are separate rings that double as the rotating element (the main disk for AZ); a press-in 3D-printed retainer holds the ball complement.

The trade-off is surface finish: printed races have higher $R_a$ than machined steel bearings. Two properties of the rotator keep this acceptable: (i) the operating speed at the output is bounded by the motor output divided by the external ratio (approximately 6~RPM on the AZ axis via the 15:40 reduction, approximately 11~RPM on EL via the 40:55 reduction, at 95\,\% duty; \S\ref{sec:slew-speed}), so contact-surface sliding velocity is low; and (ii) the bearing preload is low because the herringbone gear pair below imposes negligible axial thrust load (\S\ref{sec:gearing}). Endurance characterisation beyond the field sessions is future work.

\begin{figure}[H]
    \centering
    \includegraphics[width=0.6\textwidth]{photos/3DP_bearings_close-up.png}
    \caption{3D-printed deep-groove ball bearing close-up with 3\,mm steel balls and press-in retainer.}
\end{figure}

\begin{figure}[H]
    \centering
    \includegraphics[width=0.7\textwidth]{photos/3DP_bearings_all.png}
    \caption{Both printed bearings assembled, with the separate AZ inner race printed flat-down.}
\end{figure}

\subsection{Gearing}
\label{sec:gearing}

All gears use a double-helical (herringbone) tooth profile with a 15\degree\ helix angle, 4/3~mm normal module, and 8~mm face width. The herringbone profile cancels the axial thrust generated by each helical half on a single-helical gear, which matters here because the printed radial deep-groove bearings tolerate much less axial than radial load. The 4/3~mm module is non-standard in the ISO 53 preferred series~\cite{iso53} but gives integer reference diameters (e.g. 15 teeth $\times$ 4/3~mm = 20~mm) that print cleanly on FDM at 0.2~mm layer height.

The gear-ratio bookkeeping:

\begin{table}[H]
    \centering
    \footnotesize
    \begin{tabular}{@{}lllllllc@{}}
        \toprule
        Axis & Teeth & Ext. ratio & Internal gearbox & Total ratio & CPR & Ticks/\degree & Travel \\
        \midrule
        AZ & 15:40 & 2.667:1 & 516:1 & 1376:1  & 64 & 244.6 & $\pm$360\degree \\
        EL & 40:55 & 1.375:1 & 516:1 & 709.5:1 & 64 & 126.1 & 0--180\degree \\
        \bottomrule
    \end{tabular}
    \caption{Rotator gear ratios, encoder counts, and per-axis positional resolution.}
\end{table}

Ticks per degree at the output is derived from the encoder count per motor revolution multiplied by the total reduction, divided by 360:

\begin{itemize}
    \item AZ: $\text{ticks/deg} = \dfrac{64 \times 516 \times (40/15)}{360} = \dfrac{64 \times 516 \times 2.667}{360} = \dfrac{88\,064}{360} = 244.6$ ticks/deg, giving a positioning resolution of $1/244.6 = 0.00409\degree$ per tick.
    \item EL: $\text{ticks/deg} = \dfrac{64 \times 516 \times (55/40)}{360} = \dfrac{64 \times 516 \times 1.375}{360} = \dfrac{45\,408}{360} = 126.1$ ticks/deg, giving 0.00793\degree\ per tick.
\end{itemize}

Gears press-fit directly onto the motor's 6~mm D-shafts with additional set-screws.

\begin{figure}[H]
    \centering
    \includegraphics[width=0.7\textwidth]{photos/3DP_gears_all.png}
    \caption{Printed herringbone gear set: AZ pair (15:40, 2.667:1) and EL pair (40:55, 1.375:1).}
\end{figure}

\begin{figure}[H]
    \centering
    \includegraphics[width=0.7\textwidth]{photos/gears_cad.png}
    \caption{Parametric CAD source of the gear profile, generated in Onshape.}
\end{figure}

\section{Motors}
\label{sec:motors}

The axis motor is the FAPG36-555-EN (FONEACC): a 24~V brushed DC motor with a 516:1 internal planetary gearbox~\cite{fapg36-555}. Manufacturer spec: approximately 16~RPM no-load at 24~V, nominal output torque 4~Nm, 6~mm D-shaft output. A Hall-effect quadrature encoder is mounted on the motor shaft (before the gearbox). The manufacturer datasheet specifies 12~PPR (48~CPR in 4$\times$ decode); the purchased units empirically measure as 16~PPR (64~CPR in 4$\times$ decode), verified by matching firmware encoder counts against known commanded axis rotations during bring-up. The 6~mm D-shaft press-fits directly into the printed gear.

\subsection{Torque budget}

Hanssens~\cite{hanssens2023} identified insufficient stepper torque as the failure mode of his azimuth rotation tests. The DC-motor path chosen here multiplies the manufacturer output torque by the external gear ratio:

\begin{itemize}
    \item AZ: $4~\text{Nm} \times 2.667 = 10.7$~Nm at the output shaft (manufacturer spec $\times$ external ratio).
    \item EL: $4~\text{Nm} \times 1.375 = 5.5$~Nm at the output shaft.
\end{itemize}

\paragraph{Measured output torque.} The assembled rotator was bench-measured with a lever-arm-plus-kitchen-scale setup. A rigid arm attached to the axis output pushes down on the scale at a known radius while the axis is commanded to rotate at 95\,\% duty; the peak scale reading gives the output force.

\begin{itemize}
    \item AZ: 0.71~kgf at 0.50~m lever $\Rightarrow T_\text{AZ} = 0.71 \cdot 9.81 \cdot 0.50 = 3.48$~Nm (Fig.~\ref{fig:torque-az}).
    \item EL: 1.10~kgf at 0.40~m lever $\Rightarrow T_\text{EL} = 1.10 \cdot 9.81 \cdot 0.40 = 4.32$~Nm (Fig.~\ref{fig:torque-el}).
\end{itemize}

Measured EL reaches approximately 80\,\% of the spec-derived 5.5~Nm. Measured AZ reaches approximately 33\,\% of the spec-derived 10.7~Nm; the AZ rig lets the rotor continue to move as the scale deflects, so the reading is the rotating-under-load torque rather than the stalled peak. Both measured values exceed what the Siretta Oscar 44 Yagi (mass $\sim$0.6~kg) requires at the lever of the assembled antenna.

\begin{figure}[H]
    \centering
    \begin{minipage}{0.48\textwidth}
        \centering
        \includegraphics[width=\textwidth]{photos/torque_az_arm.png}
    \end{minipage}\hfill
    \begin{minipage}{0.48\textwidth}
        \centering
        \includegraphics[width=\textwidth]{photos/torque_az_scale.png}
    \end{minipage}
    \caption{AZ torque measurement: 0.50~m lever, scale reads 0.71~kgf $\Rightarrow$ 3.48~Nm.}
    \label{fig:torque-az}
\end{figure}

\begin{figure}[H]
    \centering
    \begin{minipage}{0.48\textwidth}
        \centering
        \includegraphics[width=\textwidth]{photos/torque_el_arm.png}
    \end{minipage}\hfill
    \begin{minipage}{0.48\textwidth}
        \centering
        \includegraphics[width=\textwidth]{photos/torque_el_scale.png}
    \end{minipage}
    \caption{EL torque measurement: 0.40~m lever, scale reads 1.10~kgf $\Rightarrow$ 4.32~Nm.}
    \label{fig:torque-el}
\end{figure}

\subsection{DC vs stepper}

The choice of DC motors over stepper motors was informed by Hanssens' torque test~\cite{hanssens2023} and by three properties of the DC-motor-plus-high-ratio-gearbox path:

\begin{enumerate}
    \item Zero idle current with mechanical position holding. A brushed DC motor draws no current when stationary, and a 516:1 reduction makes back-driving from the output side require $516\times$ the static friction torque of the unloaded motor rotor, which for a planetary gearbox of this ratio practically prevents back-drive under wind gusts. The station is unpowered at the motor terminals for the $\sim$90\,\% of the orbit it is idle (about 80 minutes out of a 92-minute orbit, per \S\ref{sec:orbital-mechanics}). In contrast, a stepper requires continuous holding current (typically 0.8--1.5~A per phase at rated current per the Pololu stepper catalogue~\cite{pololu-stepper}) to resist back-drive; a de-energised stepper has no holding torque.
    \item Torque retained at low speed. The 516:1 internal gearbox plus the external stage delivers the torque budget listed above across the full output-speed range used for tracking (0--35\degree/s). Steppers lose torque above their pull-out speed and require microstepping to produce smooth motion, which further reduces available torque per step.
    \item Closed-loop position verification. The quadrature encoder on the motor shaft gives continuous position feedback. A stepper skipping a step (wind gust, obstruction, or torque-curve excursion) silently diverges commanded from actual position unless an external encoder is added; with the encoder already in place here, the PID loop sees any discrepancy within one control period.
\end{enumerate}

The trade-off is firmware complexity: the DC-motor path requires a PID loop and quadrature-decode ISRs (\S\ref{sec:rotator-pid}, \S\ref{sec:quadrature}). The encoder feedback also makes safety features possible (runaway detection, soft-limit enforcement) that a stepper-without-encoder cannot implement.

No slip ring is fitted on the azimuth axis. The firmware tracks cumulative rotation and enforces soft limits at $\pm$360\degree, so the antenna cable never wraps more than one full turn. Between passes, the \code{PARK} command returns the antenna to true-north / zero position via the shortest path, which unwraps any accumulated wrap before the next observation (\S\ref{sec:easycomm}).

\section{Motor Controller PCB}

The v1 MotorPCB is a 4-layer KiCad~9 design with 13 hierarchical sub-sheets, JLCPCB assembled from LCSC basic parts.

Components:
\begin{itemize}
    \item RP2040 (Pi Pico module): EasyComm firmware, PID, encoder ISRs.
    \item 2$\times$ TB6642FG~\cite{tb6642fg}: 50~V absolute-max / 4.5~A peak H-bridge. IN1/IN2 + PWM. 20~kHz PWM frequency.
    \item ADXL345~\cite{adxl345}: 3-axis accelerometer on SPI0, gravity vector for EL homing.
    \item Encoder inputs: 4 GPIO with IRQ, 4$\times$ quadrature decoding. RC filter caps removed (see \S\ref{sec:bringup}).
    \item NeoPixel: optional WS2812B status output: idle = green, tracking = blue, on-target = cyan, fault = red.
\end{itemize}

The fielded v1 firmware uses GP16 for AZ IN1 after the GP15 damage described in \S\ref{sec:bringup}; the schematic's original GP15 assignment is retained only as bring-up history.

\begin{table}[H]
    \centering
    \begin{tabular}{lll}
        \toprule
        Function & AZ pins & EL pins \\
        \midrule
        Motor IN1/IN2/PWM & GP16/GP14/GP22 & GP28/GP27/GP21 \\
        Encoder A/B       & GP11/GP10      & GP12/GP13 \\
        \bottomrule
    \end{tabular}
    \caption{Rotator PCB pin assignments (as-deployed v1).}
\end{table}

\begin{figure}[H]
    \centering
    \includegraphics[width=0.7\textwidth]{photos/AZEL_control_front.png}
    \caption{KiCad 3D render of the v1 MotorPCB.}
\end{figure}

\begin{figure}[H]
    \centering
    \includegraphics[page=1,width=\textwidth]{schematics/MotorPCB_schematic.pdf}
    \caption{Top-level hierarchical schematic of the v1 MotorPCB. The 13 sub-sheets (RP2040 controller, solder-pad breakouts, encoder front-end, AZ/EL H-bridges, IMU, status LED, USB, power, etc.) are exported in full in the project KiCad repository (\file{MotorPCB/}).}
    \label{fig:motorpcb-schematic}
\end{figure}

\section{Bring-Up and Hardware Debugging}
\label{sec:bringup}

The first-revision MotorPCB shipped with three defects, all fixed with bodge wires and reflected in the v2 design:
\begin{itemize}
    \item Missing EL driver traces. OUT1/OUT2 on the elevation TB6642FG were unrouted; bodged through unused NC pins (6, 11).
    \item Encoder RC filters (C1--C4). 100~nF caps on the Hall-sensor lines low-passed the signal into unusability; removed. Debouncing is handled in firmware.
    \item AZ driver blown by bootloader cycling. The 1200-baud bootloader trick with 24~V connected damaged the TB6642FG IN1 ESD diode and the GP15 pad; driver replaced and IN1 rerouted to GP16.
\end{itemize}

The v2 MOTOR\_RF\_HAT corrects all three natively: routed EL outputs, no encoder caps, and 100\,$\Omega$ series resistors (R11--R14) on every GPIO-to-driver line. Merging the motor controller and RF front-end onto one Pi HAT also removes the USB cable and separate 5\,V path.

\section{Cost}
\label{sec:cost}

Two hardware configurations are priced separately. The first is the as-deployed v1 setup that was used for field validation (a separate MotorPCB and RF~HAT, connected to the Pi via USB and UART respectively); v1 existed to validate the architecture and firmware and to let the RF and rotator subsystems be developed in parallel. The second is the v2 combined MOTOR\_RF\_HAT, the thesis's final hardware deliverable: design-complete, routed, BOM finalised, ready to order but not yet fabricated at the time of writing. v2 reaches $\sim$€40 for the combined controller/RF PCB, $\sim$€75 for the rotator hardware excluding the Pi, and $\sim$€100 for the station core, which misses the strict ``under €50 full rotator'' target once motors and mechanics are included but lands at approximately one-eighth of the cheapest commercial unit (€760 Yaesu G-5500~\cite{yaesu-g5500}).

Component prices from the LCSC BOM tool (April 2026), bare PCB from JLCPCB (5-piece order). Full BOM in Appendix~\ref{app:bom} and \file{ThesisPaper/bom\_complete.csv}.

\subsection{v1 (deployed, field-validated)}

\begin{figure}[H]
    \centering
    \includegraphics[width=0.7\textwidth]{photos/rf_hat_bench.jpg}
    \caption{v1 RF HAT: 6-layer Pi HAT containing the dual-CC1200 transceivers.}
    \label{fig:rf-hat-front}
\end{figure}

\begin{figure}[H]
    \centering
    \includegraphics[width=0.7\textwidth]{photos/motorcontrol.jpg}
    \caption{v1 MotorPCB assembled prototype, mounted in the rotator electronics bay.}
\end{figure}

\begin{table}[H]
    \centering
    \begin{tabular}{lr}
        \toprule
        Component & \euro{} \\
        \midrule
        MotorPCB (KiCad 9, 4-layer, 13 sub-sheets, JLCPCB assembled) & $\sim$30 \\
        RF HAT (6-layer Pi HAT, 2$\times$ CC1200, JLCPCB assembled) & $\sim$55 \\
        Raspberry Pi 3 Model A+ & 25 \\
        Mechanical (motors, ASA, bearings, hardware) & $\sim$35 \\
        \midrule
        \textbf{v1 total} & \textbf{$\sim$€145} \\
        \bottomrule
    \end{tabular}
    \caption{v1 (as-deployed) cost breakdown.}
\end{table}

\subsection{v2 (MOTOR\_RF\_HAT, design-complete, ready to order)}

The v2 board merges motor control and RF front-end onto a single 6-layer Pi HAT. The A4950 (SOIC-8-EP) replaces the TB6642FG (HSOP-16) for availability (the TB6642FG went out of stock at JLCPCB during the project), and four 100\,$\Omega$ series resistors (R11--R14) on the GPIO-to-driver lines limit transients during bootloader cycling. One RP2040 handles both motor control and RF, so the USB cable and separate 5~V path are removed. Hand assembly avoids JLCPCB's extended-part setup fee ($\sim$€18/board for the many precise RF L/C values).

\begin{figure}[H]
    \centering
    \includegraphics[width=0.7\textwidth]{photos/rf_hat_motor_front.png}
    \caption{v2 MOTOR\_RF\_HAT (front): 6-layer Pi HAT merging motor control and dual-CC1200 RF.}
    \label{fig:motor-rf-hat-front}
\end{figure}

\begin{figure}[H]
    \centering
    \includegraphics[width=0.7\textwidth]{photos/rf_hat_motor_back.png}
    \caption{v2 MOTOR\_RF\_HAT (back).}
\end{figure}

\begin{table}[H]
    \centering
    \begin{tabular}{lr}
        \toprule
        Component & \euro{} \\
        \midrule
        \multicolumn{2}{l}{\textbf{MOTOR\_RF\_HAT PCB}} \\
        Components (LCSC, incl. MOQ rounding) & 28.44 \\
        SMA connectors + 15~nH inductors (not on LCSC) & 2.00 \\
        Raspberry Pi Pico (hand-soldered) & 4.00 \\
        Bare PCB (6-layer, JLCPCB, per board from 5-pc order) & 6.00 \\
        \textbf{PCB subtotal} & \textbf{$\sim$40} \\
        \midrule
        \multicolumn{2}{l}{\textbf{Off-board}} \\
        FAPG36-555-EN motors (2$\times$) & 16.00 \\
        ASA filament ($\sim$400~g) & 10.00 \\
        Steel balls, screws, heat-set inserts & 5.00 \\
        PVC pipe + misc hardware & 2.00 \\
        ADXL345 IMU breakout & 1.50 \\
        \textbf{Mechanical subtotal} & \textbf{$\sim$35} \\
        \midrule
        \multicolumn{2}{l}{\textbf{Station computer}} \\
        Raspberry Pi 3 Model A+ & 25.00 \\
        \midrule
        \textbf{v2 station-core total} & \textbf{$\sim$€100} \\
        \bottomrule
    \end{tabular}
    \caption{v2 (MOTOR\_RF\_HAT, projected) cost breakdown. Excludes antenna, coax, power supply, and tripod.}
\end{table}

Excludes antenna, coax, power supply, and tripod (user-supplied). v2 brings the rotator BOM (PCB + motors + mechanics, excluding the Pi) to $\sim$€75, or $\sim$€40 for the combined controller/RF PCB alone. This misses a strict full-rotator €50 target but stays well below existing open-source and commercial rotators (Table~\ref{tab:rotator-comparison}).
